\documentclass[12pt]{article}

\usepackage[a4paper, margin=1in]{geometry}
\usepackage{amsmath,amssymb,amsthm}
\usepackage{mathtools}
\usepackage{mathrsfs}
\usepackage{array}
\usepackage{booktabs}
\usepackage{tabularx}
\usepackage{enumitem}
\usepackage{float}
\usepackage{xcolor}
\usepackage[hidelinks]{hyperref}
\usepackage{fontspec}

\defaultfontfeatures{Ligatures=TeX}
\IfFontExistsTF{Times New Roman}
  {\setmainfont{Times New Roman}}
  {\setmainfont{TeX Gyre Termes}}
\IfFontExistsTF{Menlo}
  {\setmonofont{Menlo}}
  {\setmonofont{Latin Modern Mono}}

\newtheorem{proposition}{Proposition}[section]
\theoremstyle{definition}
\newtheorem{definition}[proposition]{Definition}
\theoremstyle{remark}
\newtheorem{remark}[proposition]{Remark}

\newcolumntype{P}[1]{>{\raggedright\arraybackslash}p{#1}}
\newcolumntype{Y}{X}

\newcommand{\AU}{\mathcal A_U}
\newcommand{\Ac}{\mathcal A_c}
\newcommand{\Bbrain}{\mathcal B_{\mathrm{brain}}}
\newcommand{\Bint}{\mathcal A_{\mathrm{int}}^{\mathrm{st}}}
\newcommand{\Bintpre}{\mathcal A_{\mathrm{int}}^{\circ}}
\newcommand{\Ash}{\mathcal A_{\mathrm{sh}}}
\newcommand{\Kern}{K(t,s)}
\newcommand{\Phiij}{\Phi_{ij}(t)}
\newcommand{\tightlist}{%
  \setlength{\itemsep}{0pt}\setlength{\parskip}{0pt}}

\title{%
  \textbf{The Brain as a Plastic Observable-Algebra Network}\\
  \textbf{Internal Shared Stable Hulls, Neural Plasticity, and First-Person Report}%
}

\author{
  Sidong Liu, PhD \\
  \small iBioStratix Ltd \\
  \small \texttt{sidongliu@hotmail.com}
}

\date{May 2026\\
\small Zenodo DOI: \href{https://doi.org/10.5281/zenodo.20044329}{10.5281/zenodo.20044329}}

\begin{document}
\emergencystretch=2em
\raggedbottom
\hbadness=10000
\vbadness=5000

\maketitle

\begin{abstract}
The self-model is often treated as a point-like subject, while the brain
is often described as an information-processing machine. This paper is a
standalone anchoring interface paper: it gives a structural interface
between an observable-algebra framework and one biological observer
regime, without extending the Foundation formal theory and without
making an empirical neuroscience claim.

The core proposal is that a brain-centered embodied observer regime can
anchor an accessible algebra \(\Ac\) at biological scale. Neural
read-out structures are read as idealized local subalgebras
\(\mathcal B_i\), synaptic and functional couplings are modeled by
normal unital completely positive maps \(\Phi_{ij}(t)\), coupling
history is represented by \(K(t,s)\), and the reportable self-model is
read as an internal shared stable hull.

All claims are conditional and interface-level. Neuroscience examples
are used only as compatible re-readings, not as evidence for the
framework. The paper does not prove qualia, does not derive an
empirically testable construction of \(\Bbrain\), and does not authorize
consciousness migration, digital immortality, or first-person identity
transfer.
\end{abstract}

\clearpage
\tableofcontents
\clearpage

\section{Introduction: The Need for a Biological Anchor}

The Genesis and Consciousness papers describe consciousness-adjacent
structure at the level of accessible algebras, shared sectors, modular
flow, source-level feedback, and self-description limits. Those papers
do not, by themselves, specify a concrete biological observer regime. A
reader can therefore ask a simple anchoring question: if an accessible
algebra \(\Ac\) is not a free-floating abstraction, how is it anchored in
the biological case?

This paper answers that question only at the level of a framework
interface. It does not introduce Foundation IV. It does not extend the
formal Foundation series. It also does not present an empirical
neuroscience model with estimated parameters, data fits, or clinical
predictions.

The proposed interface is that the brain-centered embodied observer
regime can be read as a biological anchor for an accessible algebra. In
that regime, the brain is a central plastic algebraic organ, but it is
not the whole observer regime. The body, sensorimotor loops,
interoception, affective systems, environmental coupling, tools,
records, and other agents also constrain what is accessible and
reportable.

The brain is not the source of consciousness in a reductive sense, nor
a container of an inner world. Consciousness, in this framework, is not
a thing that requires a source, but a structural modality of an
accessible algebra anchored through a biological observer regime. The
brain is such a central anchoring component: it helps stabilize a
plastic internal shared hull that can later be reported as ``I.''

The guiding contrast is therefore not between a ghostly subject and a
physical brain. The guiding contrast is between a many-channel,
plastic, history-dependent observer regime and the later reportable
self-model that this regime stabilizes. The reportable ``I'' is not
discarded as unreal. It is read as an effective stable structure within
\(\Ac\), not as an independent point-subject.

\subsection*{Contribution}

The paper makes three contributions. First, it gives a self-contained
biological anchoring interface for the observable-algebra framework. A
reader need not accept any empirical claim about the brain to understand
the interface: the key question is how an accessible algebra can be
anchored through an embodied observer regime.

Second, the paper turns the internal self-model into a formal target.
Instead of treating the ``I'' as either a primitive subject or a mere
illusion, it asks what kind of stable readable substructure could be
supported by multiple internal read-out structures. The answer is
conditional: under a selected coarse-graining, specified conditional
expectations, and specified coupling maps, the candidate overlap
generates an internal shared stable hull
(Propositions~\ref{prop:hull-construction-main}--\ref{prop:common-private-main}).

Third, the paper separates the algebraic contribution from the
neuroscience interface. Neuroscience examples are used to show where the
language might touch existing phenomena, not to validate the framework.
This separation is essential. Without it, the paper would risk borrowing
neuroscience authority for a formal structure that has not been derived
from data. With it, the paper can be read as an algebra-first interface
proposal whose empirical relevance remains a later research program.

\subsection*{How to read the structural claims}

Most claims in the paper have the form ``\(X\) can be read as \(Y\)''
rather than ``\(X=Y\).'' This is not rhetorical caution alone. It marks
the level of the paper. The biological phenomenon is not being reduced
to an algebraic object, and the algebraic object is not being inferred
from the biological phenomenon. The paper proposes a partial structural
reading between the two.

This also explains why the paper uses conditional propositions rather
than unconditional theorems. The propositions are not claims that brains
automatically satisfy the hypotheses. They say that, if an embodied
observer regime is idealized by the specified operator-algebraic datum,
then the internal hull and common/private decomposition are well-defined
within that datum. The formal burden is therefore explicit rather than
hidden.

\subsection*{What this paper does not do}

\noindent\fbox{%
\begin{minipage}{0.94\linewidth}
\begin{itemize}[leftmargin=1.5em]
\item It does not solve the hard problem.
\item It does not prove that the brain produces consciousness.
\item It does not propose that consciousness has a source requiring identification.
\item It does not derive an empirically testable construction of \(\Bbrain\) from neuroscience data.
\item It does not replace Global Workspace Theory, predictive processing, higher-order theories, or clinical neuroscience.
\item It does not authorize consciousness migration, digital immortality, or first-person identity transfer.
\end{itemize}
\end{minipage}}

\subsection*{Plan of the paper}

Section~2 gives the minimal self-contained observable-algebra
vocabulary needed for the paper. Section~3 explains why accessible
algebras require embodied observer regimes. Section~4 formulates the
brain-level reading as a plastic observable-algebra network. Section~5
describes the self-model as an internal shared stable hull.

Section~6 gives the formal core: a conditional hull construction, a
conditional common/private decomposition, and a fragmentation criterion.
Section~7 gives a limited compatible re-reading of first-person report
timing. Section~8 compresses the neuroscience interfaces into a table
and selected discussions. Section~9 separates guardrails from
limitations. Section~10 states a research program, and Section~11
restates the core anchor claim. Appendix~A gives the fuller formal
skeleton.

\section{A Minimal Self-Contained Framework: Imported Vocabulary}

This section gives the imported minimal vocabulary required for the
present paper. The terms are compatible with Foundation~I, Foundation~II, and
Foundation~III of the Consciousness Series~\cite{FoundationI,FoundationII,FoundationIII},
but the definitions below are intended to be readable without consulting
those papers. Nothing in this section is meant to re-prove the
Foundation results.

\begin{definition}[Observable algebra]
An observable algebra \(\mathcal A\) is an algebraic language for
quantities that can be read, stabilized, transformed, or reported by a
specified regime. In this paper \(\mathcal A\) is idealized as a von
Neumann algebra. The von Neumann setting is used because it makes states,
normal maps, conditional expectations, and weak operator closure
available in one common language. This is a structural idealization, not
a claim that neural tissue literally presents a naturally given von
Neumann algebra.
\end{definition}

The word ``observable'' should be read structurally rather than
psychologically. An observable is not necessarily an item that is
consciously noticed. It is an element of the formal language through
which a regime can, in principle, stabilize distinctions, update states,
or support later report. This distinction matters because the paper is
not trying to identify conscious contents one by one with algebra
elements. It is specifying the structural level at which access,
readability, and report can be discussed.

The von Neumann idealization is deliberately stronger than what the
neuroscience interface can presently justify. It is used because the
formal questions of this paper require closure, states, and conditional
expectations. A weaker \(C^*\)-algebraic language would be possible in
principle, but the present construction repeatedly refers to normal
states and \(\sigma\)-weak closure. For that reason, the von Neumann
setting is the cleanest operator-algebraic envelope for a first
standalone interface paper.

\begin{definition}[Ambient algebra]
The ambient algebra \(\AU\) denotes the background observable algebra
relative to which accessible substructures are selected. It is not
treated here as a collection of already separated objects. It is the
formal substrate from which finite regimes induce accessible sectors.
\end{definition}

The ambient algebra is not an empirically specified object in this
paper. It is a bookkeeping device for the fact that access is partial.
The biological observer regime does not encounter an already segmented
world and then add subjective labels to it. Rather, the formal picture
starts from a broader observable context and asks which substructure is
available under the constraints of a finite regime.

\begin{definition}[Accessible algebra]
An accessible algebra \(\Ac\subseteq\AU\) is the observable algebra
available to a finite observer regime. It is not an arbitrary
psychological viewpoint. It is the algebraic expression of what can be
read, stabilized, updated, reported, and acted through under the
constraints of that regime.
\end{definition}

This is the main self-contained hinge of the paper. An accessible
algebra is not a mind, a brain region, or a representation stored in a
head. It is the formal expression of a regime's accessible observables.
For the biological case, this means that \(\Ac\) is anchored through the
brain-centered embodied regime, but it is not identical with the brain.
The brain is a central anchoring component, while \(\Ac\) is the
algebraic object whose access conditions are constrained by the whole
regime.

\begin{definition}[Observer regime]
An observer regime is the physical, biological, and historical set of
conditions through which an accessible algebra is anchored. An
accessible algebra only has operational meaning relative to such a
regime. In the biological case considered here, the regime is
brain-centered but embodied: body, sensorimotor loops, interoception,
environmental coupling, tools, records, and social interaction all
constrain access.
\end{definition}

The observer regime is the paper's substitute for a point-like subject.
It is also the paper's guard against a brain-in-a-vat reading. A finite
observer regime has sensors, effectors, metabolic constraints, bodily
states, memory traces, and environmental couplings. These conditions do
not merely deliver inputs to a detached algebra. They help determine
which observables can be stabilized and which distinctions can become
reportable.

This regime language is also meant to avoid a common ambiguity in
consciousness theory. A theory may speak as if there is a subject to
whom contents appear, and then ask how the brain generates those
contents. The present framework reverses that order. It first specifies
the conditions under which an accessible algebra is available. Only
then does it ask how a reportable self-model can stabilize inside that
accessible algebra. The subject is not inserted as a primitive term at
the beginning.

The approach therefore does not deny first-person structure. It denies
that first-person structure should be treated as an unanalysed point.
The reportable first person is a late and effective structure within a
regime. It has practical and ethical force because the regime can plan,
remember, repair, and be trained. But its force does not require it to
be an independent owner of the algebra.

\begin{definition}[Shared stable hull]
A shared stable hull is a stable overlap sector that can be read across
multiple source chains or regimes under specified compatibility
conditions. Foundation~II develops this notion for external
multi-source settings~\cite{FoundationII}. The present paper uses the
idea by mechanism-homology only: it asks whether a structurally similar
overlap can be read inside one accessible algebra.
\end{definition}

The external shared hull and the internal shared hull should not be
confused. The external construction concerns overlap across source
chains or observer regimes. The internal construction developed here
concerns overlap among read-out structures inside one accessible
algebra. The formal analogy is useful because both cases ask for a
stable readable sector under compatibility conditions. It is limited
because the index set, coupling maps, and interpretation of private
complement are different.

\begin{definition}[Private complement]
The private complement is the part of a regime or read-out channel that
does not enter the shared-readable sector. In external settings it marks
what does not become public overlap. In the internal setting considered
here, it marks what remains local to a bodily, affective, memory, or
report-limited channel.
\end{definition}

The private complement is not a hidden metaphysical substance. It is a
formal way of saying that not everything shaping a reportable self-model
is itself shared-readable across the relevant channels. A bodily state
may bias report without becoming linguistic. An affective salience
signal may constrain action without becoming a public reason. The
private complement records this asymmetry without treating it as a
separate subject.

\begin{definition}[History-dependent coupling]
The symbol \(K(t,s)\) denotes a history-dependent coupling object. It is
not a single neural kernel derived from data. It is a structural
placeholder for non-Markovian dependence: past interactions change which
future states are easier to read, synchronize, report, or act through.
\end{definition}

The notation \(K(t,s)\) is intentionally generic. It can be read as a
history-dependent kernel, a coupling-history object, or a record of how
prior interactions constrain later synchronization. At the biological
interface, synaptic plasticity, network reconfiguration, habit
formation, and affective conditioning may all provide partial empirical
interfaces to this role. The paper does not identify \(K(t,s)\) with
any one of them.

The same caution applies to \(\Phi_{ij}(t)\). A coupling map is not a
wire carrying a pre-existing message from one component to another. It
is a structural representation of how one read-out structure can change
the future readability or stabilization conditions of another. In
neural language, such coupling may be mediated by synapses, recurrent
loops, neuromodulation, functional connectivity, or body-mediated
feedback. In the present paper these are not collapsed into one
biophysical mechanism. They are placed under one formal heading only at
the level of structural readability.

The point of the present paper is not to identify the brain with
\(\Ac\). The point is to describe how a brain-centered embodied regime
can anchor one biological instance of an accessible algebra. In that
reading, \(\Ac\) remains the algebraic object, while the biological
regime is the anchoring condition.

\begin{figure}[H]
\centering
\fbox{%
\begin{minipage}{0.9\linewidth}
\centering
\[
\AU
\longrightarrow
\Ac
\longrightarrow
\{\mathcal B_i\}_{i\in I}
\longrightarrow
\Bint
\]
\[
\text{\small ambient algebra}
\quad
\text{\small accessible algebra}
\quad
\text{\small read-out structures}
\quad
\text{\small internal stable hull}
\]
\end{minipage}}
\caption{Schematic only. The figure shows the minimal framework used in
this paper: \(\AU\), \(\Ac\), internal read-out structures
\(\{\mathcal B_i\}\), and the internal shared stable hull \(\Bint\).}
\end{figure}

\subsection{Translation Conventions}

Because the paper is an interface paper, several terms have two
readings: an operator-algebraic reading and a biological or
phenomenological interface reading. The following table fixes the
translation conventions used throughout the paper.

\begin{table}[H]
\centering
\small
\begin{tabularx}{0.98\linewidth}{P{0.24\linewidth}P{0.35\linewidth}Y}
\toprule
\textbf{Term} & \textbf{Formal role} & \textbf{Interface role} \\
\midrule
\(\AU\) &
Ambient observable algebra. &
Background structural context, not an empirically measured universe. \\
\(\Ac\) &
Accessible von Neumann subalgebra. &
What a finite embodied regime can read, stabilize, update, and report. \\
\(\mathcal B_i\) &
Idealized read-out subalgebra at a fixed scale. &
Neural, bodily, memory, affective, or linguistic channel read as a local structure. \\
\(E_i\) &
Faithful normal conditional expectation onto \(\mathcal B_i\). &
Idealized selection of channel-readable content. \\
\(\Phi_{ij}(t)\) &
Normal unital completely positive coupling map. &
Plastic coupling between read-out structures. \\
\(K(t,s)\) &
History-dependent coupling object. &
Memory, habit, plasticity, or prior coupling history at structural level. \\
\(\Bint\) &
\(\sigma\)-weakly closed internal hull candidate. &
Stable shared-readable self-model, later reportable as ``I.'' \\
\bottomrule
\end{tabularx}
\caption{Translation conventions. The biological column gives
compatible interface readings only; it does not identify the biological
phenomenon with the formal object.}
\end{table}

The table is also a warning against category mistakes. A cortical
network is not literally a von Neumann algebra; it is read through one
when the paper asks about observable access and reportability. A
synapse is not literally a completely positive map; synaptic and
functional coupling can be modeled by such maps at a selected structural
level. The value of the formalism lies in making the conditions of
shared readability explicit, not in replacing biology with algebra.

\subsection{Why This Minimal Vocabulary Is Enough}

The paper does not require the full machinery of the Foundation series.
It uses only four structural ideas. First, access is partial and
regime-dependent. Second, stable overlap can be studied algebraically.
Third, not everything local to a channel becomes shared-readable.
Fourth, history changes future readability. These four ideas are enough
to formulate the brain-level anchor without importing the full source
chain, modular-flow, or perfuming-loop machinery.

This is why the paper can be standalone. Foundation~I--III provide the
background motivation and a broader research program, but the argument
below only needs the self-contained vocabulary introduced here. A reader
who accepts the definitions in this section can follow the later
construction even without accepting the whole Consciousness Series.

\section{Accessible Algebras Require Embodied Observer Regimes}

An accessible algebra is specified relative to an observer regime. It is
not a free-floating abstraction. What can be read, updated, reported,
and acted through depends on the embodied conditions under which
accessibility is realized.

For the biological case, the observer regime is brain-centered but not
brain-only. The brain is the central plastic algebraic organ of the
regime. The full regime also includes the body, sensorimotor loops,
interoceptive and affective systems, environmental coupling, tools,
records, habits, and social interaction.

This distinction prevents a common drift. The paper does not claim that
the brain is \(\Ac\). It claims that the brain-centered embodied regime
anchors an accessible algebra. The accessible algebra is constrained by
the whole embodied regime, even when the brain provides the most
plastic and integrative component.

\begin{figure}[H]
\centering
\fbox{%
\begin{minipage}{0.9\linewidth}
\begin{center}
\textbf{Brain-centered embodied observer regime}
\end{center}
\[
\begin{array}{c}
\text{brain: central plastic algebraic organ}\\
\downarrow\\
\text{body and sensorimotor loops}\\
\downarrow\\
\text{interoception and affective subsystems}\\
\downarrow\\
\text{environmental coupling: light, sound, gravity, other agents}\\
\downarrow\\
\text{long-term inscriptions: habits, tools, records}
\end{array}
\]
\end{minipage}}
\caption{Schematic only. The brain is central in the biological anchor,
but it is not the whole observer regime.}
\end{figure}

The embodied regime also prevents a second drift, namely the
brain-in-a-vat reading. The accessible algebra is not anchored by an
isolated neural object floating apart from sensorimotor and
environmental constraints. It is anchored by a coupled regime whose
neural component is central but not sufficient by itself.

This paper therefore treats the brain as an anchoring component rather
than as a container of a world. The world as read by the regime is not
inside the brain. The readable conditions under which a world is
stabilized are constrained through the regime.

This regime-level language also explains why the paper does not start
from an isolated computational brain. A purely neural description may be
appropriate for many empirical purposes, but the accessible algebra of a
biological observer is constrained by more than neural tissue. Posture,
interoceptive load, sensorimotor contingencies, sleep history, tools,
written records, and social interaction can change what is readable and
reportable. These factors do not have to be modeled in detail here, but
they justify treating the brain as central within an embodied regime
rather than as the whole regime.

The point is not to dilute the role of the brain. The brain remains the
dominant plastic integration organ of the biological anchor. The point
is to avoid making the algebraic object \(\Ac\) identical to a neural
object before the regime of access has been specified. In the notation
used later, a candidate brain algebra \(\Bbrain\) may be a crucial
substructure of \(\Ac\), but the embodied regime also constrains the
state, read-out maps, report windows, and environmental couplings.

\section{The Brain as a Plastic Observable-Algebra Network}

The brain can be read, at a chosen biological coarse-graining scale, as
a plastic observable-algebra network. This is a structural reading, not
an empirical derivation of a neural algebra from data. The reading is
useful because it allows neural plasticity, coupling, memory, and
self-model stabilization to be discussed in the same language.

Two overstatements must be ruled out at the start. A neural read-out
unit is not a quantum state. A synapse is not literal quantum
entanglement. Both claims are too strong, and both would push the paper
toward a quantum-brain hypothesis that is not being made.

Instead, choose a biological coarse-graining scale. At that scale, a
read-out unit may be a single neuron, a neural ensemble, a cortical
column, a functional network, a sensorimotor loop, a memory subsystem,
an affective subsystem, or a language subsystem. We denote the
corresponding effective read-out subalgebras by
\(\{\mathcal B_i\}_{i\in I}\), with \(\mathcal B_i\subseteq\Ac\).

At the biological interface, a neural read-out unit is not a passive
node but a local integrator. Dendritic inputs, excitatory and inhibitory
modulation, and thresholded action-potential generation already
implement a local read-out operation before network-level coupling is
considered~\cite{StuartSpruston2015}.
Network-level coupling \(\Phi_{ij}(t)\) then composes such local
integrations into larger patterns.

The objects \(\{\mathcal B_i\}_{i\in I}\) are called effective
read-out structures in the main text. In the idealized setting of
Appendix~A they are assumed to be von Neumann subalgebras of \(\Ac\).
At the biological interface this idealization is a modeling choice, not
a claim that every neural ensemble is literally closed under
multiplication, adjoint, and operator-topology limits.

The choice of coarse-graining is not a cosmetic detail. At a single
neuron scale, one might emphasize dendritic integration, spike
thresholds, and local input-output selectivity. At an ensemble scale,
one might emphasize synchrony, oscillatory coordination, and functional
connectivity. At a behavioral subsystem scale, one might emphasize
sensorimotor affordances, affective salience, memory retrieval, or
language report. The notation \(\mathcal B_i\) is deliberately flexible
enough to cover these possibilities, but the formal construction in
Section~6 always fixes one scale before making a statement.

This prevents a common false precision. The paper does not claim that
there are exactly \(n\) biologically privileged subalgebras, nor that a
single neuron, cortical column, and language subsystem all belong to the
same clean algebraic level. It says that, once a scale is selected, the
relevant read-out structures can be idealized as subalgebras for the
purpose of asking whether a stable internal overlap can be constructed.

\begin{table}[H]
\centering
\small
\begin{tabularx}{0.98\linewidth}{P{0.22\linewidth}P{0.33\linewidth}Y}
\toprule
\textbf{Scale} & \textbf{Possible read-out structure} & \textbf{Reason for caution} \\
\midrule
Cellular &
Single neuron or local dendritic integration. &
Too fine-grained for report-level self-modeling by itself. \\
Mesoscopic &
Neural ensemble, column, or recurrent local circuit. &
May support stable patterns, but boundaries are model-dependent. \\
Network &
Functional network or large-scale integration pattern. &
Closer to report and action, but farther from cellular mechanism. \\
Embodied subsystem &
Sensorimotor loop, affective subsystem, memory system, language system. &
Useful for self-model analysis, but mathematically idealized as a subalgebra. \\
\bottomrule
\end{tabularx}
\caption{Possible coarse-graining scales for \(\mathcal B_i\). The
formal construction fixes one scale; it does not assert invariance
across all scales.}
\end{table}

\begin{table}[H]
\centering
\begin{tabularx}{0.96\linewidth}{P{0.34\linewidth}Y}
\toprule
\textbf{Biological structure} & \textbf{Structural reading} \\
\midrule
Neural read-out unit \(i\) &
Local effective read-out subalgebra \(\mathcal B_i\) at a chosen biological coarse-graining scale. \\
Synaptic or functional coupling \(ij\) &
Plastic coupling channel or algebraic edge \(\Phi_{ij}(t)\). \\
Plasticity &
History-dependent modification of coupling channels. \\
Memory &
Coupling history, not a static storage file. \\
Learning &
Modification of future readability, synchronization, and action-candidate stabilization. \\
Self-model &
Network-stabilized internal shared hull. \\
\bottomrule
\end{tabularx}
\caption{A structural reading, not an equality table. Each entry is to be
read as ``can be modeled as'' at a chosen coarse-graining scale.}
\end{table}

The key shift is that a coupling is not read as an information wire. The
wire metaphor is useful in engineering contexts, but it suggests that
pre-existing information objects travel along edges. The algebraic-edge
reading is different: a coupling channel changes which future states can
be jointly stabilized and read.

At the neural interface, this means that synaptic coupling is better
read as a plastic electrochemical interface than as a fixed information
line. Chemical synapses convert electrical activity into transmitter
release and then back into postsynaptic electrical change; the
structural point is not the biochemical detail itself, but the fact that
coupling is state-dependent, modulatory, and plastic~\cite{Sudhof2012,Sudhof2013}.
This supports the use of \(\Phi_{ij}(t)\) as
a coupling-channel notation rather than a wire notation.

In this sense, a synaptic or functional coupling is a plastic edge in a
network of local read-out algebras. Each experience modifies the future
conditions of readability. It changes which channels synchronize, which
states become easier to stabilize, and which action candidates become
more likely to enter reportable dynamics.

\paragraph{Why not a standard neural network model?}
Standard neural network models can describe input-output mappings,
learned weights, prediction errors, and task performance. The present
paper is addressing a different structural layer: the algebraic
conditions of observable access, the conditions under which states are
readable, the distinction between shared and private sectors, the
location of the report boundary, and the role of history-dependent
coupling. This is not a competing neural network model. Rather, a neural
network model can be read as one possible concrete coupling instance
inside a broader \(\Phi_{ij}(t)\)-style framework, if the relevant
coarse-graining and read-out maps are independently specified.

\begin{figure}[H]
\centering
\fbox{%
\begin{minipage}{0.9\linewidth}
\[
\mathcal B_1
\xleftrightarrow{\ \Phi_{12}(t)\ }
\mathcal B_2
\xleftrightarrow{\ \Phi_{23}(t)\ }
\mathcal B_3
\quad \cdots \quad
\mathcal B_n
\]
\[
\text{plastic coupling history } K(t,s)
\quad \leadsto \quad
\text{future synchronization and stable read-out}
\]
\end{minipage}}
\caption{Schematic only. The figure represents a plastic
observable-algebra network. It is not a connectome, a data plot, or a
claim of measured neural coupling.}
\end{figure}

This reading is compatible with ordinary neuroscience vocabulary, but
it is not a replacement for it. It does not specify neurotransmitters,
spike timing, synaptic weights, or network physiology. It gives a
structural interface in which such details, if independently modeled,
could potentially constrain a future construction of \(\Bbrain\).

\begin{remark}[Candidate brain algebra]
Here \(\Bbrain\) denotes a full candidate observable algebra for the
brain component within the embodied observer regime. If such a
\(\Bbrain\) exists at a chosen coarse-graining scale, it is understood
to satisfy \(\Bbrain\subseteq\Ac\) and \(\Bbrain\neq\Ac\): the brain
component anchors, but does not exhaust, the accessible algebra of the
embodied regime. The paper does not construct \(\Bbrain\) from
neuroscience data, and it does not assume that such a construction is
unique.
\end{remark}

\section{Internal Shared Stable Hulls and the Self-Model}

The self-model is read here as an internal shared stable hull within a
single accessible algebra. It is not a new independent algebra outside
\(\Ac\). It is a stable internal substructure produced by overlap,
coupling, and recurrent calibration among multiple read-out
subalgebras.

The analogy with Foundation~II is partial. Foundation~II constructs an
external shared stable hull across source chains under named
conditions~\cite{FoundationII}. Here the internal reading concerns
subalgebras within one embodied accessible algebra. The two settings are
mechanism-homologous, not strictly isomorphic.

The contrast can be stated more explicitly. In the external case, the
problem is how multiple source chains or regimes can share a stable
overlap while retaining private complements. In the internal case, the
problem is how multiple read-out structures inside one embodied
accessible algebra can support a stable reportable self-model while
retaining local bodily, affective, memory, or pre-report dynamics. The
same words ``shared'' and ``private'' are used, but the index sets and
interpretation differ.

This is why the paper does not simply copy the external construction.
The internal construction requires its own datum: read-out subalgebras
\(\mathcal B_i\), conditional expectations \(E_i\), coupling maps
\(\Phi_{ij}(t)\), a state or effective state \(\omega\), and a
history-dependent coupling object \(K\). These objects identify the
chosen biological coarse-graining scale. Without such a datum,
``internal hull'' would remain only a metaphor.

At the internal scale, different read-out subalgebras contribute
different aspects of the reportable self-model. Sensorimotor channels
stabilize bodily boundary and action affordance. Interoceptive and
affective channels stabilize urgency, salience, and valence. Memory
channels stabilize temporal continuity. Language channels stabilize
reportable narrative.

The reportable self-model is not simply identical with any one of these
channels. It is the stable overlap generated by their recurrent
coupling. When the overlap is strong and coordinated, the regime reports
``I am doing this'' or ``I decided this.'' When the overlap weakens or
fragments, the reportable structure may become less unified.

This stable overlap is what the paper calls the internal shared stable
hull \(\Bint\) at the informal level. Appendix~A gives a provisional
formalization of this object.

The construction route is simple but conditional. First choose a scale
at which read-out structures can be idealized as subalgebras. Then
choose conditional expectations that select channel-readable content.
Then identify candidate overlap observables by comparing their
state-readings across relevant channels. Finally take the unital
\(*\)-algebra generated by those candidates and close it in the
\(\sigma\)-weak topology. Section~6 states this route as a conditional
construction rather than as a biological theorem.

This route also clarifies what the self-model is not. It is not a point
that owns the channels. It is not the language channel alone. It is not
the entire accessible algebra. It is a stable readable substructure
within \(\Ac\), dependent on the selected scale, coupling family, state,
and report window.

This does not mean that the self-model is merely an illusion. The
self-model is effective within \(\Ac\). It can plan, report, repair,
remember, and assume responsibility. What is denied is not its
effectiveness, but its status as an independent point-subject.

The internal hull also has a private side. Some bodily, affective, or
pre-report dynamics may shape the hull without becoming fully
shared-readable across internal channels or externally reportable to
other observers. This is why the paper distinguishes first-person
report from the full content of experience.

The relation to Paper~I is motivational rather than central. Paper~I
argues that geometry and modular flow are not independently deformable
in the relevant algebraic setting~\cite{PaperI}. The present paper uses
that lesson only as a caution against separating ``objective brain'' and
``subjective report'' into two independently specified domains.

\section{Formal Core}

\paragraph{Mathematical status of the results.}
The formal contribution of this section is a pair of conditional
construction statements (Propositions~\ref{prop:hull-construction-main}
and~\ref{prop:common-private-main}) together with a fragmentation
criterion. These are conditional constructions within standard von
Neumann algebra theory, not new existence results in operator algebra
or algebraic quantum field theory. Their value is structural: they show
that the informal phrase ``internal shared stable hull'' admits a
well-defined formal reading once a suitable operator-algebraic datum is
specified. The paper does not claim a new von Neumann algebra existence
theorem, a new classification result, or a strengthening of the
Tomita--Takesaki apparatus. It claims that the interface between the
observable-algebra framework and a biological observer regime can be
made precise at the level described below.

This section states the formal core in a compact form. Appendix~A gives
the fuller skeleton and the limitations of the construction. The
results below are conditional: they assume a fixed coarse-graining
scale, specified conditional expectations, specified coupling maps, and
a selected report window.

The purpose of the section is modest but precise. It does not construct
a brain algebra from data and it does not prove that biological brains
satisfy the hypotheses. It shows that, once a biological interface has
been idealized by a suitable operator-algebraic datum, the phrase
``internal shared stable hull'' can be given a concrete formal reading.
The contribution is therefore not an unconditional theorem about
consciousness. It is a conditional bridge from the structural language
of the paper to a well-defined operator-algebraic object.

\subsection{Setup}

An internal read-out datum is a tuple
\[
\mathfrak D_{\mathrm{int}}
=
(\Ac,\omega,\{\mathcal B_i\}_{i\in I},\{E_i\}_{i\in I},
\{\Phi_{ij}(t)\}_{i,j\in I},K),
\]
where \(\Ac\) is an accessible von Neumann algebra, \(\omega\) is a
faithful normal state or specified effective state, \(\mathcal B_i\) are
idealized read-out von Neumann subalgebras of \(\Ac\),
\(E_i:\Ac\to\mathcal B_i\) are faithful normal conditional
expectations, \(\Phi_{ij}(t):\Ac\to\Ac\) are normal unital completely
positive maps modeling coupling channels, and \(K\) records
history-dependent modulation of the coupling family.

The datum is not derived from neuroscience data in this paper. It is a
formal interface object. Its role is to make explicit which analytic
choices would be needed before one could speak about an internal stable
hull at a fixed biological coarse-graining scale.

Each component of the datum records a different idealization decision.
The family \(\{\mathcal B_i\}\) records the chosen biological grain. The
maps \(E_i\) record how channel-readable content is selected. The maps
\(\Phi_{ij}(t)\) record how coupling is represented at that grain. The
state \(\omega\) records which expectations or effective readings are
being compared. The kernel \(K\) records the history-dependence of the
coupling family. Changing any of these choices may change the resulting
candidate hull.

\subsection{Conditional Hull Construction}

Let
\[
\mathcal O_{\mathrm{ov}}
=
\bigl\{\,a\in\Ac :
\omega(E_i(a))=\omega(E_j(a))
\text{ for all relevant }i,j\in I
\bigr\},
\]
with equality understood up to the tolerance specified by the regime.
This set is not intended to contain all experiential content. It is the
candidate overlap locus for those observables whose channel-selected
readings agree across the relevant subfamilies. The tolerance language
is necessary because the biological interpretation is finite-regime and
coarse-grained, not exact in the sense of an ideal mathematical
measurement.

Assume the following conditions:
\begin{enumerate}[label=(H\arabic*)]
\item[(H0)] The maps \(E_i:\Ac\to\mathcal B_i\) are faithful normal
conditional expectations onto von Neumann subalgebras
\(\mathcal B_i\subseteq\Ac\).
\item \textbf{Coupling stability.} The normal unital completely positive
maps \(\Phi_{ij}(t)\), restricted to the unital \(*\)-algebra generated
by \(\mathcal O_{\mathrm{ov}}\), preserve the generated candidate hull
up to the regime tolerance over the relevant report window.
\item \textbf{State compatibility.} The restricted states
\(\omega|_{\mathcal B_i}\) agree on \(\mathcal O_{\mathrm{ov}}\) up to
the regime tolerance, and the overlap is non-trivial in the selected
regime.
\item \textbf{Conditional-expectation compatibility.} The maps \(E_i\)
and the coupling family \(\{\Phi_{ij}(t)\}\) are compatible on
\(\mathcal O_{\mathrm{ov}}\) up to the regime tolerance.
\end{enumerate}

\paragraph{Interpretation of the hypotheses.}
Condition (H0) is the strongest mathematical idealization. It says that
the selected read-out structures are not merely informal channels but
von Neumann subalgebras admitting faithful normal conditional
expectations. This is not guaranteed in general and is not claimed to
follow from neuroscience. It is the assumption that lets the paper
compare channel-selected observables inside one algebraic environment.

Condition (H1) is the stability assumption. It is stated at the level of
the generated candidate hull because normal unital completely positive
maps do not generally preserve products. Stability cannot be inferred
from pointwise preservation of a generating set. It must be required for
the candidate hull over the selected report window. This is why the
proposition is a conditional construction rather than an existence
theorem for all data.

Condition (H2) prevents the overlap from collapsing into a merely formal
object. If the only shared-readable elements were scalar multiples of
the identity, the construction would be formally possible but
phenomenologically empty. The paper therefore treats non-trivial overlap
as an explicit assumption. Later work could try to identify empirical or
formal constraints under which such non-triviality is plausible.

Condition (H3) ties the channel selection maps to the coupling maps. It
requires that coupling and read-out do not immediately destroy the
candidate overlap. In biological language, this corresponds to the idea
that the structures supporting a reportable self-model must remain
coordinated long enough to be read, attributed, or acted through. The
condition is finite-window stability, not permanent identity.

\begin{table}[H]
\centering
\small
\begin{tabularx}{0.98\linewidth}{P{0.16\linewidth}P{0.38\linewidth}Y}
\toprule
\textbf{Hypothesis} & \textbf{Formal burden} & \textbf{Interface meaning} \\
\midrule
(H0) &
Existence of faithful normal conditional expectations onto selected subalgebras. &
Channel-readable content can be selected at the chosen scale. \\
(H1) &
Generated candidate hull is stable under restricted coupling maps over a report window. &
The self-model candidate persists long enough to be read or attributed. \\
(H2) &
Restricted states agree non-trivially on the overlap locus. &
The overlap is not merely scalar or empty at the selected scale. \\
(H3) &
Conditional expectations and coupling maps are compatible on the overlap locus. &
Read-out and coupling do not immediately destroy the shared candidate. \\
\bottomrule
\end{tabularx}
\caption{Interpretation of the formal hypotheses. The right column is
only an interface reading, not an empirical derivation of the
hypotheses.}
\end{table}

\begin{proposition}[Conditional Hull Construction]
\label{prop:hull-construction-main}
Fix an internal read-out datum \(\mathfrak D_{\mathrm{int}}\) satisfying
(H0)--(H3). Then \(\mathcal O_{\mathrm{ov}}\) generates a unital
\(*\)-algebra
\[
\Bintpre\subseteq\Ac.
\]
Its \(\sigma\)-weak closure \(\Bint\) is a von Neumann subalgebra of
\(\Ac\). Under (H1), the generated candidate hull is invariant under the
restricted action of the coupling family \(\{\Phi_{ij}(t)\}\) over the
relevant report window, up to the regime tolerance.
\end{proposition}

\paragraph{Proof sketch.}
By (H0), \(\mathcal O_{\mathrm{ov}}\) is a well-defined subset of
\(\Ac\) selected by the state \(\omega\) and the conditional
expectations \(E_i\). Any subset of a unital \(*\)-algebra generates a
unital \(*\)-algebra, giving \(\Bintpre\subseteq\Ac\). Its
\(\sigma\)-weak closure in the von Neumann algebra \(\Ac\) is a von
Neumann subalgebra. Non-triviality is part of (H2), not a derived
consequence. The stability statement is exactly (H1), which is stated at
the level of the generated candidate hull rather than merely at the
level of the generating set.

The proposition is intentionally conditional in two ways. First, it
requires the existence of channel expectations \(E_i\), which is not
automatic. Second, it requires a tolerance-stable overlap that is not
merely scalar. These assumptions identify the formal burden of any later
empirical or mathematical strengthening. If they fail, the framework
does not say that a unified internal hull exists at that scale.

\subsection{Conditional Common/Private Read-Out Decomposition}

\begin{proposition}[Conditional Common/Private Decomposition]
\label{prop:common-private-main}
Suppose, in addition to Proposition~\ref{prop:hull-construction-main},
that there exists a faithful normal conditional expectation
\[
E_{\mathrm{int}}:\Ac\to\Bint.
\]
For each read-out subalgebra \(\mathcal B_i\), define
\[
\mathcal B_i^{\mathrm{com}}
=
\mathcal B_i\cap\Bint
\]
and let
\[
\mathcal B_i^{\mathrm{priv}}
=
\ker(E_{\mathrm{int}}|_{\mathcal B_i})
\]
in the sense of the restricted map. Then
\(\mathcal B_i^{\mathrm{com}}\) is the shared-readable part of
\(\mathcal B_i\) relative to the internal hull, while
\(\mathcal B_i^{\mathrm{priv}}\) captures what remains local to channel
\(i\) under the selected conditional expectation.
\end{proposition}

\paragraph{Proof sketch.}
The intersection \(\mathcal B_i\cap\Bint\) lies in both the channel
subalgebra and the internal hull, so it is the natural shared-readable
candidate. The restricted kernel of \(E_{\mathrm{int}}\) contains those
elements of \(\mathcal B_i\) that vanish under the selected read-out
onto \(\Bint\). This gives a conditional common/private decomposition
relative to \(E_{\mathrm{int}}\), the state \(\omega\), the
coarse-graining scale, and the report window. It is not claimed to be
unique and is not claimed to be an algebraic direct-sum decomposition in
general.

The notation \(\mathcal B_i^{\mathrm{priv}}\) should therefore be read
with care. It does not mean that the private part is invisible to the
whole organism or causally inert. It means that, relative to the chosen
map \(E_{\mathrm{int}}\), those elements do not contribute to the
selected shared-readable hull. They may still influence behavior,
affect, attention, or later report through other channels.

\begin{remark}[Additional assumption]
The existence of \(E_{\mathrm{int}}\) is an extra assumption. A faithful
normal conditional expectation from \(\Ac\) onto \(\Bint\) is not
automatic in a general von Neumann algebra. This is the conditional
content of Proposition~\ref{prop:common-private-main}.
\end{remark}

\paragraph{Why this decomposition matters.}
The common/private distinction is not added for mathematical ornament.
It is the formal place where the paper separates reportable self-model
from the full range of processes shaping experience and behavior. A
memory trace, bodily tension, affective bias, or subreport motor
tendency may influence the self-model without being shared-readable in
the selected hull. Conversely, what enters
\(\mathcal B_i^{\mathrm{com}}\) is available to the internal overlap
relative to the chosen \(E_{\mathrm{int}}\).

This gives a precise sense in which the self-model can be effective
without being exhaustive. It can guide report and action because it is a
stable shared-readable structure. It can fail to capture everything
because each channel may retain a private complement. The framework
therefore avoids both extremes: it does not treat the self-model as a
fiction, and it does not treat the self-model as the whole mind.

\subsection{Failure Modes}

The same formal setup also identifies ways in which the construction may
fail. These failure modes are important because they prevent the paper
from treating a unified self-model as automatic.

\begin{enumerate}[label=(M\arabic*)]
\item The selected read-out structures may not admit the required
conditional expectations. In that case the chosen scale is not suitable
for this construction.
\item The overlap locus \(\mathcal O_{\mathrm{ov}}\) may be trivial. In
that case the generated hull may exist formally but fail to support a
meaningful shared-readable self-model.
\item The coupling maps may fail to preserve the generated candidate
hull over the report window. In that case the self-model candidate is
too unstable for the intended reading.
\item The common/private decomposition may fail if no suitable
\(E_{\mathrm{int}}\) exists. In that case the hull construction may
still be available, but the private complement analysis remains
under-specified.
\end{enumerate}

These failure modes are not pathologies in the clinical sense. They are
formal warnings. A future empirical or mathematical application would
need to show not only that a datum can be written down, but that the
chosen datum avoids these failure modes in the regime under study.

\subsection{Fragmentation Criterion}

The internal hull may fragment when the conditions supporting stable
overlap fail. A formal fragmentation regime may be indicated by one or
more of the following:
\begin{enumerate}[label=(F\arabic*)]
\item The coupling graph on \(I\) disconnects into weakly coupled or
uncoupled components.
\item The overlap among relevant subfamilies falls below a threshold
set by the regime.
\item Conditional-expectation compatibility fails across relevant
subfamilies.
\item State compatibility fails across relevant subfamilies over the
report window.
\end{enumerate}

Under such conditions, multiple local self-model candidates may become
structurally possible. This is a formal criterion, not a diagnosis. It
does not identify DID, split-brain, or any clinical state.

\section{First-Person Report and Timing}

First-person report is treated as a late, reportable phase of a broader
internal stabilization process. This section uses Libet-style timing
only as a compatible re-reading. It is not evidence for the framework,
and it is not a general theory of decision.

In standard presentations of Libet-style wrist-flexion experiments, a
readiness potential is reported before the subject's reported awareness
of deciding, and both precede the action~\cite{Libet1983,Haggard2008,Schurger2012}. The timing values usually associated
with this paradigm are approximate and paradigm-dependent. They are
drawn from self-paced, low-stakes actions, and extrapolation to complex
or consequential decisions is not supported by the present framework.

The readiness potential is treated here as a window into
action-preparation dynamics, not as a direct marker of completed
decision. In the present language, it can be read as an early measurable
trace of action-candidate stabilization. It does not show that the
internal shared hull model is true.

\begin{table}[H]
\centering
\begin{tabularx}{0.95\linewidth}{P{0.18\linewidth}Y P{0.22\linewidth}}
\toprule
\textbf{Timing} & \textbf{Compatible structural re-reading} & \textbf{Observable marker} \\
\midrule
Approx. \(-550\) ms &
Action-candidate stabilization begins to become measurable in the relevant motor-preparation regime. &
Readiness potential. \\
Approx. \(-200\) ms &
Reportable self-attribution becomes available: ``I decided.'' &
W time report. \\
Action time &
Motor output occurs. &
Movement. \\
\bottomrule
\end{tabularx}
\caption{Schematic timing diagram. It is not an empirical data plot and
does not generalize to complex or consequential decisions.}
\end{table}

The compatible re-reading is that the reportable ``I decided'' is not
the first event in the action process. It belongs to the stage at which
an action tendency has become stable enough to be attributed to the
self-model. This does not make first-person report worthless. It
locates it in the reportable phase of a longer internal process.

This also blocks a misleading opposition. It is not ``the brain decided
and I was bypassed.'' The internal hull is stabilized within the
person-level regime through which action is prepared, attributed, and
executed. What is misleading is the identification of the person with
the late linguistic report alone.

Responsibility is therefore not erased. The relevant responsible
structure is the embodied regime with its history-dependent couplings,
plasticity, repair capacities, and future modifiability. Responsibility
is not located in a single sentence, ``I decided,'' but in a trainable
and accountable condition-chain.

\section{Neuroscience Interfaces}

This section gives interface points only. It is not a review article,
and it does not claim to replace the theories or clinical categories it
mentions. The cases below are selective examples rather than a complete
inventory. For all empirical phenomena discussed below, the phenomenon
is used as compatible re-reading only, not as evidence for the
framework. The neuroscience citations in this section are orientation
markers: they locate where the algebraic language might touch existing
empirical or theoretical vocabularies. They are not evidential premises
for the formal results in Section~6, and they do not make the framework
empirically confirmed. External references in this section require
author verification before submission.

The table is included to prevent a misleading impression. The paper is
not collecting neuroscientific evidence for the algebraic framework.
Instead, it identifies where existing empirical or theoretical
vocabularies would touch the algebraic objects introduced above:
coupling maps, report boundaries, external calibration, private
complements, and hull fragmentation. A reader should therefore treat
the table as an interface map, not as an argument from neuroscience to
the framework.

\begin{table}[H]
\centering
\small
\begin{tabularx}{0.98\linewidth}{P{0.18\linewidth}P{0.34\linewidth}Y}
\toprule
\textbf{Interface} & \textbf{External theory or phenomenon} & \textbf{Structural reading in this framework} \\
\midrule
Split-brain &
Split-brain cases involve disrupted interhemispheric communication~\cite{Sperry1968,Gazzaniga2000}. &
Can be read as coupling disconnection and possible internal-hull fragmentation. \\
DID &
DID is a clinically complex dissociative condition and should not be diagnosed from a formal model~\cite{APA2022,ISSTD2011}. &
Can be read as protected decoupling among memory, affective, bodily, and narrative channels. \\
Sleep &
Sleep involves reduced sensory processing and motor output with memory and network maintenance~\cite{RaschBorn2013,TononiCirelli2014}. &
Can be read as offline reweighting of \(K(t,s)\) under reduced immediate action demand. \\
Dreaming &
Dreaming involves reduced external input with internally active memory, affective, and imagery systems~\cite{NirTononi2010,Hobson2009}. &
Can be read as weakened external calibration while internal-hull stabilization continues. \\
Meditation &
Meditation research studies training of attention, interoception, and affective regulation~\cite{TangHolzelPosner2015,LutzDunneDavidson2007}. &
Can be read as modulation of attention and the report boundary. \\
OBE / altered self-location &
OBE and dissociative states involve altered body-ownership and self-location~\cite{Blanke2002,BlankeArzy2005}. &
Can be read as body-anchoring and self-location decoupling or remapping. \\
GWT &
Global Workspace Theory describes consciousness in terms of global availability~\cite{Baars1988,DehaeneChangeux2011}. &
Can be read as a reportable phase of internal-hull stabilization. \\
Predictive processing &
Predictive processing describes perception and action through prediction and error dynamics~\cite{Friston2010,Clark2013}. &
Can be read as one class of coupling update within the embodied observable-algebra network. \\
HOT &
Higher-order theories locate awareness in states about other states~\cite{Rosenthal2005,LauRosenthal2011}. &
Can be read as one read-out structure making another available at the report boundary. \\
\bottomrule
\end{tabularx}
\caption{Selective neuroscience and consciousness-theory interfaces.
Each structural reading is a compatible re-reading only, not evidence
for the framework and not a replacement for the external theory.}
\end{table}

The table also explains why the section does not expand every row into a
separate literature review. The most important rows for the formal core
are split-brain, global availability, and sleep/dreaming. Split-brain
interfaces with fragmentation, global availability interfaces with
reportable stabilization, and sleep/dreaming interfaces with
history-dependent coupling and external calibration. The remaining rows
are included to show where the same language may connect later, not to
claim coverage of clinical or contemplative phenomena.
None of the structural readings in the table replaces the external
theory or clinical category it interfaces with.

\paragraph{Split-brain and fragmentation.}
Split-brain cases provide the cleanest interface to the formal
fragmentation criterion in Section~6. When a major coupling route is
disrupted, some read-out subfamilies may no longer support one unified
candidate hull. This does not mean that the formal model diagnoses
split-brain patients. It only says that coupling disconnection can be
read as one biological interface to hull fragmentation.

\paragraph{Global availability and report.}
Global Workspace Theory is the most direct theory-level interface
because it focuses on global availability across multiple systems. In
the present framework, global availability can be read as a reportable
phase of internal-hull stabilization across read-out structures. The
paper does not replace GWT, does not derive its architecture, and does
not claim that GWT is wrong.

\paragraph{Sleep, dreaming, and history-dependent coupling.}
Sleep and dreaming provide the most direct interface to \(K(t,s)\) and
external calibration. Sleep can be read as a state in which immediate
external read-out and action demand are reduced while coupling history
is reweighted; dreaming can be read as a state of altered external
calibration in which transient world-model and self-model structure may
continue without direct external anchoring. The paper does not offer a
complete theory of sleep or dreaming and does not treat dream content as
evidence for the framework.

\paragraph{First-person timing as a restricted interface.}
The Libet-style discussion in Section~7 is included because it gives a
restricted timing case for first-person report. Readiness potential is
read as a measurable trace of action-preparation dynamics, not as a
direct marker of completed decision. This interface is limited to
self-paced, low-stakes action paradigms and is not generalized to
complex or consequential decisions.

\section{Guardrails and Limitations}

\subsection{Guardrails}

\paragraph{1. Not a quantum-brain hypothesis.}
This paper is not a quantum-brain hypothesis. The algebraic reading is
allowed to operate at an effective or classical-observable level, and it
does not require long-lived quantum coherence in the brain. The
counter-claim not made is: ``the brain is a quantum computer that
produces consciousness.''

\paragraph{2. Not an empirical derivation of a testable construction of \(\Bbrain\).}
This paper is not an empirical derivation of a testable construction of
\(\Bbrain\). It introduces a structural interface that future work could
potentially constrain using neuroscience data. The counter-claim not
made is: ``neuroscience data already determine the observable algebra of
the brain.''

\paragraph{3. Not exhaustive over psychological functions.}
This paper is not exhaustive over psychological functions, including
qualia, emotion, imagery, clinical symptoms, or all cognition. It
describes one structural reading of internal stabilization and
first-person report. The counter-claim not made is: ``all mental
functions are explained by the plastic observable-algebra network.''

\paragraph{4. Not a proof of qualia or a solution to the hard problem.}
This paper is not a proof of qualia or a solution to the hard problem.
It does not claim to solve the hard problem of consciousness, and it
does not reduce first-person experience to third-person description. The
counter-claim not made is: ``the hard problem has been solved by neural
algebra.''

\paragraph{5. Not authorization for consciousness migration.}
This paper is not authorization for consciousness migration, digital
immortality, or first-person transfer. Functional copying, if discussed
elsewhere, does not establish identity transfer. The counter-claim not
made is: ``a copy of the network is the same first-person subject.''

\paragraph{6. Not authorization for brain-computer-interface mind reading.}
This paper does not claim that brain-computer interfaces can read,
write, or transfer first-person consciousness. BCI results can be read
as partial access to selected sensory, motor, or language-output
channels, not as direct access to the full internal hull. The
counter-claim not made is: ``recording enough neural activity would
read the mind or capture the subject.''

\paragraph{7. Not a claim about the soul, near-death experience, or post-mortem continuity.}
This paper does not prove, disprove, or redescribe any metaphysical
doctrine of the soul. Phenomena such as out-of-body experience are
treated here only as altered self-location readings within the
framework, not as evidence for or against first-person persistence
beyond the embodied regime. The counter-claim not made is: ``the soul
exists as a separable entity the framework endorses.''

\subsection{Limitations}

First, the internal hull candidate depends on the selected
coarse-graining scale and on the choice of coupling family. Different
choices may produce different internal hull candidates. The paper does
not resolve this scale- and choice-dependence.

Second, cross-scale consistency remains coarse. Synaptic plasticity may
operate over milliseconds, seconds, hours, days, and longer timescales,
while first-person report timing may be measured in a different regime.
This paper does not provide a full multi-scale neural dynamics model.

Third, the physical and neurobiological plausibility of assumptions
(H1)--(H3) in Section~6 and Appendix~A requires separate work. The assumptions are
structural placeholders. They are not derived from neurophysiology in
this paper.

Fourth, the relation between one embodied observer regime and
multi-observer external shared hulls from Foundation~II is not covered
here. The paper is about the internal anchor of one regime. It does not
construct a full bridge from internal hulls to social or intersubjective
worlds.

\section{Research Program}

The framework suggests a conservative research program. It does not
state that existing data validate the model. It asks how future data
could potentially constrain pieces of the proposed interface.

Connectome data could potentially constrain candidate coupling
topologies \(\Phi_{ij}(t)\). Such data would not by themselves identify
\(\Bbrain\), but they may provide empirical pressure on which
coarse-grained read-out units are plausible. Any such use would require
independent modeling choices.

EEG and fMRI signatures could potentially constrain large-scale
integration and fragmentation regimes. They would not directly measure
an internal shared hull. They might identify regimes in which a
simplified coupling model becomes more or less plausible.

Perturbation data could be used to test simplified cases of hull
stability. Stimulation, lesion, or pharmacological studies may create
natural interface cases. This paper does not propose a protocol and
does not treat such data as confirmation.

Plasticity models could provide coarse constraints on \(K(t,s)\)-like
history dependence. The relevant question would be how prior coupling
history changes future readability and synchronization. This remains a
research program rather than a result.

Long-term potentiation and related synaptic plasticity mechanisms could
potentially provide one empirical interface for this question~\cite{BlissLomo1973,BlissCollingridge1993}.
Such mechanisms are
compatible with reading memory as coupling history, but the present
framework does not claim that any single plasticity mechanism implements
\(K(t,s)\). They would at most constrain simplified biological cases.

\paragraph{Formal direction.}
The most direct formal next step is to replace the tolerance placeholders
in Section~6 with specified seminorms, state distances, or operational
equivalence relations. This would turn the phrase ``up to the regime
tolerance'' into a mathematical parameter. One could then ask how the
candidate hull varies as that parameter changes, and whether
fragmentation can be characterized by a threshold phenomenon.

A second formal direction is to study uniqueness and functoriality under
coarse-graining. The present paper explicitly allows different choices
of \(\{\mathcal B_i\}\), \(E_i\), and \(\Phi_{ij}(t)\) to produce
different candidate hulls. Later work could ask whether some hull
candidates are stable under refinement or coarsening of the biological
scale. Such a result would be much stronger than anything claimed here.

\paragraph{Empirical direction.}
The empirical direction is not to look for \(\Bint\) directly. The
internal hull is not an observable brain signal. A more conservative
program would seek constraints on the ingredients of the datum:
candidate read-out families, coupling topologies, perturbation
sensitivity, and history-dependent plasticity. If these ingredients can
be bounded in simplified settings, the formal hypotheses (H0)--(H3)
could become less arbitrary.

This still would not validate the full framework. It would only show
that some idealized datum is plausible for some limited biological
regime. That is the right scale of empirical ambition for an anchoring
interface paper.

\section{Conclusion}

The brain-centered embodied observer regime can be read as a biological
anchor for an accessible algebra. Within that regime, a plastic
observable-algebra network can stabilize an internal shared hull, later
reported as ``I.'' This does not solve consciousness, reduce selfhood to
neural machinery, or prove qualia. It gives a constrained structural
interface between the observable-algebra framework and brain-level
organization.

\appendix

\section{Optional Formal Skeleton for Internal Shared Stable Hulls}

This appendix gives a provisional formal skeleton. It is not a
completed Foundation-level theory. It introduces definitions and
conditional constructions that may guide later work.
Section~6 gives the main formal results in accessible form; this
appendix provides the fuller skeleton for readers who want the detailed
construction and its limitations.

The appendix repeats some material from Section~6 on purpose. The main
text presents the formal results in a compact way so that the
interdisciplinary argument can proceed. The appendix makes the
idealizations more explicit: the selected scale, the conditional
expectations, the coupling maps, the tolerance structure, and the report
window are all part of the construction. None of these is inferred from
neuroscience data in the present paper.

\subsection{Provisional Definition: Internal Read-Out Datum}

\begin{definition}[Provisional internal read-out datum]
An internal read-out datum is a tuple
\[
\mathfrak D_{\mathrm{int}}
=
(\Ac,\omega,\{\mathcal B_i\}_{i\in I},\{E_i\}_{i\in I},\{\Phi_{ij}(t)\}_{i,j\in I},K),
\]
where:
\begin{enumerate}[label=(\roman*)]
\item \(\Ac\) is an accessible von Neumann algebra associated with a
fixed embodied observer regime.
\item \(\omega\) is a faithful normal state, or a specified effective
state, on \(\Ac\).
\item \(\mathcal B_i\subseteq\Ac\) are effective read-out subalgebras
(von Neumann subalgebras) at a fixed biological coarse-graining scale.
\item \(E_i:\Ac\to\mathcal B_i\) are faithful normal conditional
expectations onto \(\mathcal B_i\). Their existence at a given
coarse-graining scale is an assumption, not automatic.
\item \(\Phi_{ij}(t):\Ac\to\Ac\) are normal unital completely positive
maps, with the understanding that \(\Phi_{ij}(t)\) models the coupling
between \(\mathcal B_i\) and \(\mathcal B_j\) at time \(t\). The action
of \(\Phi_{ij}(t)\) on overlap observables between \(\mathcal B_i\) and
\(\mathcal B_j\) is the structurally relevant part; the extension to
\(\Ac\) is a modeling convenience.
\item \(K\) is a history-dependent memory kernel or coupling-history
object controlling the time dependence of the family
\(\{\Phi_{ij}(t)\}\).
\end{enumerate}
\end{definition}

\begin{remark}[Scale-dependence]
The cardinality and identity of \(\{\mathcal B_i\}_{i\in I}\) and the
existence of the maps \(\{E_i\}\) depend on the coarse-graining scale
chosen within the embodied regime.
Proposition~\ref{prop:conditional-construction} is understood as a
conditional statement at a fixed such scale, not as an invariance claim
across scales.
\end{remark}

\begin{remark}[Why the \(E_i\) are listed as part of the datum]
In a general von Neumann algebra \(\Ac\), a faithful normal conditional
expectation \(E_i:\Ac\to\mathcal B_i\) is not guaranteed to exist. The
existence of such maps is therefore listed as part of the datum
\(\mathfrak D_{\mathrm{int}}\), not as an automatic consequence. This
makes the selection of candidate overlap observables in A.2 well
posed.
\end{remark}

\begin{remark}[Why completely positive maps are used]
The coupling maps \(\Phi_{ij}(t)\) are not assumed to be
\(*\)-homomorphisms. At the biological interface, coupling can be
modulatory, noisy, lossy, and state-dependent. Normal unital completely
positive maps give a flexible operator-algebraic language for such
effective transformations. The price of this flexibility is that
invariance cannot be inferred merely from invariance of a generating
set; it must be assumed at the level of the generated candidate hull, as
in (H1).
\end{remark}

\begin{remark}[Tolerance structure]
The tolerance appearing in (H1)--(H3) is part of the finite-regime
idealization. Biological read-out is not treated as exact equality of
all observables across all times. The tolerance records the report
window and the coarse-graining at which two channel-readings are counted
as compatible. A future formal paper would need to replace this
placeholder by a specified seminorm, state-distance, or operational
equivalence relation.
\end{remark}

\subsection{Conditional Construction Schema for an Internal Stable Hull}

The following is a conditional construction schema. It is not a general
existence theorem for all brains, all coarse-grainings, or all observer
regimes. It is written as a schema so that operator-algebraic readers
can see where the idealization decisions sit, and where later work
would have to commit to specific analytic choices.
Proposition~\ref{prop:conditional-construction} below is the detailed
schema version of Proposition~\ref{prop:hull-construction-main} in
Section~6. The two are not independent results; the appendix version
makes the idealization decisions more explicit.

\paragraph{Hypotheses.}
Fix an internal read-out datum \(\mathfrak D_{\mathrm{int}}\). Assume:
\begin{enumerate}[label=(H\arabic*)]
\item[(H0)] \textbf{Coarse-graining maps exist.} The conditional
expectations \(E_i:\Ac\to\mathcal B_i\) listed in
\(\mathfrak D_{\mathrm{int}}\) are faithful and normal, and each
\(\mathcal B_i\) is a von Neumann subalgebra of \(\Ac\).
\item \textbf{Coupling stability.} The normal unital completely positive
maps \(\Phi_{ij}(t)\), restricted to the unital \(*\)-algebra generated
by \(\mathcal O_{\mathrm{ov}}\), preserve the generated candidate hull
up to the regime tolerance over the relevant report window.
\item \textbf{State compatibility.} The restricted states
\(\omega|_{\mathcal B_i}\) agree on the candidate overlap observables
up to the tolerance specified by the regime.
\item \textbf{Conditional-expectation compatibility.} The maps \(E_i\) and the
coupling family \(\{\Phi_{ij}(t)\}\) are compatible on the candidate
overlap observables: applying \(E_i\) after \(\Phi_{ij}(t)\), and
applying \(\Phi_{ij}(t)\) before \(E_j\), agree on the candidate locus
up to the regime tolerance.
\end{enumerate}

\paragraph{Selection of candidate overlap observables.}
Define
\[
\mathcal O_{\mathrm{ov}}
=
\bigl\{\,a\in\Ac :
\omega(E_i(a)) = \omega(E_j(a))
\text{ for all relevant }i,j\in I
\bigr\},
\]
where equality is understood up to the tolerance specified by the
regime. The set \(\mathcal O_{\mathrm{ov}}\) is defined directly from
the state \(\omega\) and the conditional expectations \(\{E_i\}\); it
does not forward-reference (H1)--(H3). Hypotheses (H2)--(H3) then
impose additional structure on this set: (H2) requires that the
restricted states agree on \(\mathcal O_{\mathrm{ov}}\) up to
tolerance, and (H3) requires that the coupling family and the
conditional expectations are compatible on \(\mathcal O_{\mathrm{ov}}\).

\begin{proposition}[Conditional construction schema for an internal stable hull]
\label{prop:conditional-construction}
Under (H0)--(H3), the set \(\mathcal O_{\mathrm{ov}}\) generates a
\(*\)-algebra
\[
\Bintpre
\subseteq \Ac
\]
that is invariant under the restricted action of the coupling family
\(\{\Phi_{ij}(t)\}\), over the relevant report window, up to the regime
tolerance. If in addition
\(\Bintpre\) is closed in the \(\sigma\)-weak (or weak operator)
topology of \(\Ac\), its closure
\(\Bint\) is a von Neumann subalgebra of \(\Ac\).
\end{proposition}

\paragraph{Proof sketch.}
By (H0), the maps \(E_i\) are faithful normal conditional expectations
onto von Neumann subalgebras \(\mathcal B_i\). The set
\(\mathcal O_{\mathrm{ov}}\) is therefore a well-defined subset of
\(\Ac\): it consists of elements on which the state \(\omega\) assigns
equal values under all relevant \(E_i\). The unital \(*\)-algebra
generated by this subset inside \(\Ac\) is well-defined.

By (H2), the restricted states \(\omega|_{\mathcal B_i}\) agree on
\(\mathcal O_{\mathrm{ov}}\) up to the regime tolerance. This ensures
that \(\mathcal O_{\mathrm{ov}}\) is non-trivial in the sense relevant
for the construction: it contains observables that are jointly readable
across the subalgebras, not merely scalar multiples of the identity.
Non-triviality in this sense is an assumption of the schema, not a
derived consequence.

By (H1), the restricted action of \(\Phi_{ij}(t)\) on the generated
candidate hull stays within that candidate hull up to the regime
tolerance over the relevant report window. The condition is not a claim
of permanent stability. It is a finite-regime stability condition,
sufficient for reportable self-model formation.

By (H3), the conditional expectations \(E_i\) and the coupling family are
compatible on \(\mathcal O_{\mathrm{ov}}\), so repeated application of
\(\Phi_{ij}(t)\) followed by \(E_i\) does not push candidate overlap
observables outside the selected locus. This is the internal analogue
of requiring that a shared sector be stable under the operations used
to isolate it.

Let \(\Bintpre\) be the \(*\)-algebra generated by
\(\mathcal O_{\mathrm{ov}}\) inside \(\Ac\). The \(\sigma\)-weak
closure of \(\Bintpre\) in \(\Ac\) always exists; when \(\Bintpre\)
is itself \(\sigma\)-weakly closed, this closure equals \(\Bintpre\)
and is a von Neumann subalgebra of \(\Ac\).
The construction is therefore conditional on a closure discipline
compatible with the standard operator-algebraic setting used in the
rest of the Consciousness Series.

The content of Proposition~\ref{prop:conditional-construction} is not
that every embodied observer regime has a unique internal hull. It is
that, at a fixed coarse-graining scale and under named coarse-graining,
coupling, state, and conditional-expectation conditions, there is a
well-defined candidate internal stable hull. This is enough for an anchoring
interface paper.

\begin{remark}[What the construction does not show]
Proposition~\ref{prop:conditional-construction} does not prove that
the hypotheses hold in biological brains. It does not give an empirical
construction of \(\Bbrain\). It also does not prove uniqueness of
\(\Bint\), cross-scale invariance, or first-person identity transfer.
The schema form is deliberate: full analytic commitments (precise
choices of \(E_i\), of the coupling-channel normalization, of the
closure topology, and of the tolerance structure) are left as targets
for later formal work, not as claims of the present paper.
\end{remark}

\begin{remark}[Relation to Proposition~\ref{prop:hull-construction-main}]
Proposition~\ref{prop:conditional-construction} is the detailed
appendix version of Proposition~\ref{prop:hull-construction-main}. The
two statements have the same mathematical content. The appendix version
spells out the selection of \(\mathcal O_{\mathrm{ov}}\), the role of
the \(E_i\), and the reason the construction remains conditional.
\end{remark}

\subsection{Conditional Common / Private Decomposition}

\begin{proposition}[Conditional common/private decomposition]
\label{prop:common-private-appendix}
Suppose that the hypotheses of
Proposition~\ref{prop:conditional-construction} hold and that there
exists a faithful normal conditional expectation
\[
E_{\mathrm{int}}:\Ac\to\Bint.
\]
For each read-out subalgebra \(\mathcal B_i\), define
\[
\mathcal B_i^{\mathrm{com}}
=
\mathcal B_i \cap \Bint
\]
and
\[
\mathcal B_i^{\mathrm{priv}}
=
\ker(E_{\mathrm{int}}|_{\mathcal B_i}).
\]
Then \(\mathcal B_i^{\mathrm{com}}\) is the shared-readable part of
\(\mathcal B_i\) relative to the internal hull, while
\(\mathcal B_i^{\mathrm{priv}}\) captures what remains local to channel
\(i\) under the selected conditional expectation.
\end{proposition}

\paragraph{Proof sketch.}
The intersection \(\mathcal B_i\cap\Bint\) is contained both in the
local read-out subalgebra and in the internal hull, so it is the natural
shared-readable candidate. The restricted kernel of \(E_{\mathrm{int}}\)
contains the elements of \(\mathcal B_i\) that vanish under the selected
read-out onto \(\Bint\). The construction is therefore conditional on
the existence and choice of \(E_{\mathrm{int}}\).

This decomposition is not claimed to be unique. It depends on the
selected coarse-graining scale, the state \(\omega\), the coupling
family, the report window, and the choice of \(E_{\mathrm{int}}\). It is
not claimed to be an algebraic direct-sum decomposition in general.

The structural reading is that the common part contributes to the
stable self-model, while the private complement may influence behavior,
affect, or report without being fully shared-readable across all
internal channels. This is only a formal target. It is not a clinical or
phenomenological classification.

\subsection{Criterion: Fragmentation Regime}

The internal hull may fragment when the conditions supporting stable
overlap fail. A formal fragmentation regime may be indicated by one or
more of the following:
\begin{enumerate}[label=(F\arabic*)]
\item The coupling graph on \(I\) disconnects into weakly coupled or
uncoupled components.
\item The overlap among relevant subfamilies falls below a threshold
set by the regime.
\item Conditional-expectation compatibility fails across relevant
subfamilies.
\item State compatibility fails across relevant subfamilies over the
report window.
\end{enumerate}

Under such conditions, multiple local self-model candidates may become
structurally possible. This is a formal criterion, not a diagnosis. It
does not identify DID, split-brain, or any clinical state.

The criterion is intentionally weak. It describes a regime in which the
single-hull reading may fail or require refinement. It does not classify
all possible fragmentation phenomena and does not replace empirical or
clinical analysis.

\subsection{Remarks on Limitations of the Skeleton}

The construction in Appendix~A is conditional, not a general existence
theorem. It is fixed-scale and does not claim invariance across
biological coarse-grainings. A different scale may produce a different
family of read-out subalgebras and a different candidate hull.

The common/private decomposition and fragmentation criterion are
conditional formal tools. They require additional assumptions about
\(E_{\mathrm{int}}\), tolerance structure, and selected report windows.
They are included to show where later formal work might go. They are
not presented as completed biological or clinical classifications.

The skeleton does not derive an empirically testable construction of
\(\Bbrain\) from neuroscience data. It does not prove qualia. It does
not support first-person identity transfer. It is a constrained formal
appendix for an anchoring interface paper.

\newpage
\addcontentsline{toc}{section}{References}

\begin{thebibliography}{99}
\sloppy
\raggedright

\bibitem{FoundationI}
S.~Liu,
\emph{Foundation I: From Source to Manifestation: A Conditional Same-Root Program for Geometry and Modular Flow},
Consciousness Series (2026),
Zenodo, DOI:~10.5281/zenodo.19627786.

\bibitem{FoundationII}
S.~Liu,
\emph{Foundation II: Shared Manifestation: A Conditional Two-Source Program for Shared Accessible Sectors},
Consciousness Series (2026),
Zenodo, DOI:~10.5281/zenodo.19641372.

\bibitem{FoundationIII}
S.~Liu,
\emph{Foundation III: The Perfuming Loop: A Conditional One-Step Source-Level Feedback Program},
Consciousness Series (2026),
Zenodo, DOI:~10.5281/zenodo.19642189.

\bibitem{PaperI}
S.~Liu,
\emph{The Inseparability of Geometry and Modular Flow: A Structural Inseparability Theorem and Its Implications for the Hard Problem},
Consciousness Trilogy, Paper~I (2026),
Zenodo, DOI:~10.5281/zenodo.19607387.

\bibitem{StuartSpruston2015}
G.~J.~Stuart and N.~Spruston,
``Dendritic integration: 60 years of progress,''
\emph{Nature Neuroscience} \textbf{18} (2015), 1713--1721,
DOI:~10.1038/nn.4157.

\bibitem{Sudhof2012}
T.~C.~S{\"u}dhof,
``Calcium control of neurotransmitter release,''
\emph{Cold Spring Harbor Perspectives in Biology} \textbf{4} (2012),
a011353,
DOI:~10.1101/cshperspect.a011353.

\bibitem{Sudhof2013}
T.~C.~S{\"u}dhof,
``Neurotransmitter release: The last millisecond in the life of a
synaptic vesicle,''
\emph{Neuron} \textbf{80} (2013), 675--690,
DOI:~10.1016/j.neuron.2013.10.022.

\bibitem{Libet1983}
B.~Libet, C.~A.~Gleason, E.~W.~Wright, and D.~K.~Pearl,
``Time of conscious intention to act in relation to onset of cerebral
activity (readiness-potential): The unconscious initiation of a freely
voluntary act,''
\emph{Brain} \textbf{106} (1983), 623--642,
DOI:~10.1093/brain/106.3.623.

\bibitem{Sperry1968}
R.~W.~Sperry,
``Hemisphere deconnection and unity in conscious awareness,''
\emph{American Psychologist} \textbf{23} (1968), 723--733,
DOI:~10.1037/h0026839.

\bibitem{Gazzaniga2000}
M.~S.~Gazzaniga,
``Cerebral specialization and interhemispheric communication: Does the
corpus callosum enable the human condition?''
\emph{Brain} \textbf{123} (2000), 1293--1326,
DOI:~10.1093/brain/123.7.1293.

\bibitem{APA2022}
American Psychiatric Association,
\emph{Diagnostic and Statistical Manual of Mental Disorders},
5th ed., text rev. (DSM-5-TR),
American Psychiatric Association Publishing, Washington, DC, 2022.

\bibitem{ISSTD2011}
International Society for the Study of Trauma and Dissociation,
``Guidelines for treating dissociative identity disorder in adults,
third revision,''
\emph{Journal of Trauma \& Dissociation} \textbf{12} (2011), 115--187,
DOI:~10.1080/15299732.2011.537247.

\bibitem{RaschBorn2013}
B.~Rasch and J.~Born,
``About sleep's role in memory,''
\emph{Physiological Reviews} \textbf{93} (2013), 681--766,
DOI:~10.1152/physrev.00032.2012.

\bibitem{TononiCirelli2014}
G.~Tononi and C.~Cirelli,
``Sleep and the price of plasticity: From synaptic and cellular
homeostasis to memory consolidation and integration,''
\emph{Neuron} \textbf{81} (2014), 12--34,
DOI:~10.1016/j.neuron.2013.12.025.

\bibitem{NirTononi2010}
Y.~Nir and G.~Tononi,
``Dreaming and the brain: From phenomenology to neurophysiology,''
\emph{Trends in Cognitive Sciences} \textbf{14} (2010), 88--100,
DOI:~10.1016/j.tics.2009.12.001.

\bibitem{Hobson2009}
J.~A.~Hobson,
``REM sleep and dreaming: Towards a theory of protoconsciousness,''
\emph{Nature Reviews Neuroscience} \textbf{10} (2009), 803--813,
DOI:~10.1038/nrn2716.

\bibitem{TangHolzelPosner2015}
Y.-Y.~Tang, B.~K.~H{\"o}lzel, and M.~I.~Posner,
``The neuroscience of mindfulness meditation,''
\emph{Nature Reviews Neuroscience} \textbf{16} (2015), 213--225,
DOI:~10.1038/nrn3916.

\bibitem{LutzDunneDavidson2007}
A.~Lutz, J.~D.~Dunne, and R.~J.~Davidson,
``Meditation and the neuroscience of consciousness: An introduction,''
in P.~D.~Zelazo, M.~Moscovitch, and E.~Thompson (eds.),
\emph{The Cambridge Handbook of Consciousness},
Cambridge University Press, Cambridge, 2007, 499--551,
DOI:~10.1017/CBO9780511816789.020.

\bibitem{Blanke2002}
O.~Blanke, S.~Ortigue, T.~Landis, and M.~Seeck,
``Stimulating illusory own-body perceptions,''
\emph{Nature} \textbf{419} (2002), 269--270,
DOI:~10.1038/419269a.

\bibitem{BlankeArzy2005}
O.~Blanke and S.~Arzy,
``The out-of-body experience: Disturbed self-processing at the
temporo-parietal junction,''
\emph{The Neuroscientist} \textbf{11} (2005), 16--24,
DOI:~10.1177/1073858404270885.

\bibitem{Baars1988}
B.~J.~Baars,
\emph{A Cognitive Theory of Consciousness},
Cambridge University Press, Cambridge, 1988.

\bibitem{DehaeneChangeux2011}
S.~Dehaene and J.-P.~Changeux,
``Experimental and theoretical approaches to conscious processing,''
\emph{Neuron} \textbf{70} (2011), 200--227,
DOI:~10.1016/j.neuron.2011.03.018.

\bibitem{Friston2010}
K.~Friston,
``The free-energy principle: A unified brain theory?''
\emph{Nature Reviews Neuroscience} \textbf{11} (2010), 127--138,
DOI:~10.1038/nrn2787.

\bibitem{Clark2013}
A.~Clark,
``Whatever next? Predictive brains, situated agents, and the future of
cognitive science,''
\emph{Behavioral and Brain Sciences} \textbf{36} (2013), 181--204,
DOI:~10.1017/S0140525X12000477.

\bibitem{Rosenthal2005}
D.~M.~Rosenthal,
\emph{Consciousness and Mind},
Oxford University Press, Oxford, 2005.

\bibitem{LauRosenthal2011}
H.~Lau and D.~Rosenthal,
``Empirical support for higher-order theories of conscious awareness,''
\emph{Trends in Cognitive Sciences} \textbf{15} (2011), 365--373,
DOI:~10.1016/j.tics.2011.05.009.

\bibitem{BlissLomo1973}
T.~V.~P.~Bliss and T.~L{\o}mo,
``Long-lasting potentiation of synaptic transmission in the dentate
area of the anaesthetized rabbit following stimulation of the perforant
path,''
\emph{The Journal of Physiology} \textbf{232} (1973), 331--356,
DOI:~10.1113/jphysiol.1973.sp010273.

\bibitem{BlissCollingridge1993}
T.~V.~P.~Bliss and G.~L.~Collingridge,
``A synaptic model of memory: Long-term potentiation in the
hippocampus,''
\emph{Nature} \textbf{361} (1993), 31--39,
DOI:~10.1038/361031a0.

\bibitem{Haggard2008}
P.~Haggard,
``Human volition: Towards a neuroscience of will,''
\emph{Nature Reviews Neuroscience} \textbf{9} (2008), 934--946,
DOI:~10.1038/nrn2497.

\bibitem{Schurger2012}
A.~Schurger, J.~D.~Sitt, and S.~Dehaene,
``An accumulator model for spontaneous neural activity prior to
self-initiated movement,''
\emph{Proceedings of the National Academy of Sciences} \textbf{109}
(2012), E2904--E2913,
DOI:~10.1073/pnas.1210467109.

\end{thebibliography}

\end{document}
